% !TEX root = ../main.tex
%==============================================================
\chapter{Triển khai hệ thống}
\label{ch:trien-khai}

Chương này mô tả cách hiện thực các khái niệm ở Chương~\ref{ch:ly-thuyet} thành
một hệ thống chạy được. Nội dung lần lượt trình bày kiến trúc tổng thể và công
nghệ sử dụng; cầu nối giữa CUDA và OpenGL cùng vòng lặp kết xuất thời gian thực;
nhân Path Tracing và cấu hình thực thi trên GPU; hệ thống vật liệu và mô hình
camera; cách quản lý cảnh và dựng thế giới trên GPU; bộ nạp cảnh theo định dạng
Mitsuba~3; cơ chế tích lũy theo thời gian kèm hiệu chỉnh gamma; và cuối cùng là
giao diện điều khiển tương tác. Phần trình bày tập trung mô tả cách tổ chức, cấu
trúc dữ liệu và luồng xử lý ở mức tổng quan, thay vì đi sâu vào mã nguồn chi tiết.

\section{Tổng quan kiến trúc hệ thống}
Hệ thống được chia thành ba tầng phối hợp với nhau trong mỗi khung hình:
\begin{itemize}[noitemsep]
  \item \textbf{Tầng host (CPU, C++):} quản lý vòng lặp ứng dụng, cấp phát bộ
        nhớ GPU, nạp cảnh, xử lý nhập liệu và điều khiển camera.
  \item \textbf{Tầng device (GPU, CUDA):} thực thi nhân Path Tracing, tính màu
        cho từng điểm ảnh và ghi kết quả trực tiếp vào bộ nhớ video (VRAM).
  \item \textbf{Tầng hiển thị (OpenGL + ImGui):} đưa ảnh từ VRAM lên màn hình và
        vẽ bảng điều khiển tương tác.
\end{itemize}
Điểm mấu chốt về hiệu năng là dữ liệu ảnh không phải đi vòng qua CPU:
CUDA ghi thẳng vào bộ đệm mà OpenGL dùng để hiển thị, nhờ cơ chế CUDA--OpenGL
Interop trình bày ở Mục~\ref{sec:interop}.

\section{Công nghệ và môi trường phát triển}
Hệ thống được xây dựng bằng C++ và CUDA, cấu hình biên dịch qua CMake.
\begin{itemize}[noitemsep]
  \item Ngôn ngữ: C++ (tầng host) và CUDA C (nhân thiết bị).
  \item CUDA Toolkit phiên bản $\geq 11.0$, biên dịch bằng \texttt{nvcc}.
  \item CMake phiên bản $\geq 3.18$.
  \item Hiển thị: OpenGL kết hợp GLFW (quản lý cửa sổ, nhập liệu).
  \item Giao diện người dùng: thư viện ImGui.
  \item Sinh số ngẫu nhiên trên GPU: thư viện \texttt{cuRAND}.
  \item Phân tích cảnh Mitsuba: thư viện \texttt{tinyxml2}.
\end{itemize}

\section{Cầu nối CUDA--OpenGL và vòng lặp kết xuất}
\label{sec:interop}
Để hiển thị thời gian thực mà không tốn chi phí sao chép qua lại giữa CPU và
GPU, hệ thống dùng cơ chế interop. Một Pixel Buffer Object (PBO)
được tạo bởi OpenGL nằm trên VRAM, sau đó đăng ký với CUDA bằng
\texttt{cudaGraphicsGLRegisterBuffer}. Từ đó, CUDA có thể ghi thẳng kết quả vào
PBO, còn OpenGL dán PBO lên một texture để hiển thị.

Trong mỗi vòng lặp khung hình, hệ thống thực hiện bốn bước:
\begin{enumerate}[noitemsep]
  \item Cập nhật giao diện ImGui và camera theo thao tác chuột/bàn phím; nếu có
        thay đổi thì đồng bộ camera xuống GPU và đặt lại bộ đếm tích lũy.
  \item Khóa (\texttt{map}) PBO, lấy con trỏ VRAM rồi gọi nhân kết xuất để tô
        màu toàn bộ điểm ảnh; sau đó nhả khóa (\texttt{unmap}).
  \item OpenGL sao chép PBO vào texture (bằng DMA phần cứng) và vẽ một hình chữ
        nhật phủ kín cửa sổ.
  \item Vẽ đè giao diện ImGui và hoán đổi bộ đệm (double buffering).
\end{enumerate}

\section{Hiện thực nhân Path Tracing trên CUDA}
Nhân kết xuất chính ánh xạ mỗi luồng GPU tới đúng một điểm ảnh thông qua chỉ số
toàn cục tính từ chỉ số khối và chỉ số luồng. Mỗi luồng nạp trạng thái sinh số
ngẫu nhiên riêng của điểm ảnh mình phụ trách, làm nhiễu nhẹ vị trí bắn tia trong
ô điểm ảnh để khử răng cưa, dựng tia tương ứng từ camera rồi gọi hàm dò đường để
tính màu.

Hàm dò đường được hiện thực theo kiểu lặp thay cho đệ quy, vì đệ quy sâu tốn ngăn
xếp và không phù hợp với GPU. Thay vì gọi lại chính nó sau mỗi va chạm, hàm duy
trì một hệ số suy giảm $\beta$ và nhân dồn nó với hệ số phản xạ của vật liệu tại
mỗi lần dội, đúng theo công thức throughput ở Chương~\ref{ch:ly-thuyet}. Vòng lặp
kết thúc trong ba trường hợp: tia thoát ra ngoài cảnh và nhận màu bầu trời (nhân
với $\beta$ để trả về đóng góp của đường), vật liệu hấp thụ hoàn toàn không tán
xạ tiếp, hoặc số lần dội chạm ngưỡng tối đa (50) để bảo đảm luôn dừng. Khi tìm
giao, cận dưới tham số tia được đặt là một giá trị dương nhỏ ($t_{\min}=0{,}001$)
nhằm tránh hiện tượng shadow acne -- tia thứ cấp tự giao lại chính bề mặt vừa
xuất phát do sai số dấu phẩy động.

Về cấu hình thực thi, nhân được khởi chạy trên một lưới hai chiều phủ toàn bộ
khung hình, mỗi khối gồm $8\times8=64$ luồng tương ứng một mảng điểm ảnh liền kề.
Cách chia khối này cân bằng giữa mức độ chiếm dụng của đa xử lý luồng và tính cục
bộ khi truy cập bộ nhớ. Vì mỗi điểm ảnh cần một dãy số ngẫu nhiên độc lập cho quá
trình lấy mẫu Monte Carlo, hệ thống dùng thư viện \texttt{cuRAND} và gieo hạt
giống riêng cho từng điểm ảnh trong một nhân khởi tạo chạy một lần trước vòng lặp;
trạng thái này được lưu lại và tái sử dụng qua các khung. Một đặc thù cần lưu ý
khi tối ưu là sự phân kỳ luồng: các tia trong cùng một nhóm luồng có thể chạm vật
liệu khác loại hoặc dừng dò đường ở số lần dội khác nhau, khiến GPU phải thực thi
tuần tự các nhánh và làm giảm hiệu suất so với trường hợp lý tưởng.

\section{Hệ thống vật liệu}
Vật liệu được biểu diễn bằng một \texttt{struct Material} gọn nhẹ gồm kiểu vật
liệu, albedo, hệ số nhám \texttt{fuzz} và chiết suất \texttt{ir}. Cách dùng
\texttt{struct} phẳng (thay vì lớp ảo) giúp truyền dữ liệu xuống GPU thuận lợi.
Một hàm tán xạ chung phân nhánh theo kiểu vật liệu để xác định tia tán xạ:
\begin{itemize}[noitemsep]
  \item \textbf{Lambertian:} hướng tán xạ là pháp tuyến cộng một vector đơn vị
        ngẫu nhiên; nếu hướng thu được suy biến (gần bằng 0) thì thay bằng pháp
        tuyến.
  \item \textbf{Metal:} phản xạ gương quanh pháp tuyến, cộng thêm nhiễu trong quả
        cầu bán kính \texttt{fuzz} để mô phỏng bề mặt nhám; chỉ giữ lại tia còn
        hướng ra ngoài bề mặt.
  \item \textbf{Dielectric:} tính góc tới, xét điều kiện phản xạ toàn phần và so
        sánh xác suất Schlick~\eqref{eq:schlick} với một số ngẫu nhiên để chọn
        giữa phản xạ và khúc xạ theo định luật Snell; nhờ vậy mỗi tia chỉ đi theo
        một nhánh mà qua nhiều mẫu vẫn tái tạo đúng tỉ lệ Fresnel.
\end{itemize}

\section{Mô hình camera}
Lớp camera tiền tính các vector khung nhìn từ vị trí đặt máy, điểm nhìn, trường
nhìn, tỉ lệ khung hình, khẩu độ và khoảng lấy nét. Với mỗi tọa độ ảnh, hàm sinh
tia lấy ngẫu nhiên một điểm trên đĩa thấu kính làm gốc tia (tạo độ sâu trường
ảnh) và hướng tia đi qua điểm lấy nét tương ứng, đồng thời gán cho tia một mốc
thời gian ngẫu nhiên trong khoảng phơi sáng $[t_0, t_1]$ để phục vụ nhòe chuyển
động. Cách sinh tia này hiện thực trực tiếp mô hình thấu kính mỏng và tích phân
theo thời gian đã trình bày ở Chương~\ref{ch:ly-thuyet}.

\section{Quản lý cảnh và xây dựng thế giới trên GPU}
Trên phía host, mỗi cảnh được mô tả bằng một danh sách cấu trúc dữ liệu gọn nhẹ,
mỗi phần tử ghi tâm, bán kính, vật liệu và thông tin chuyển động của một khối cầu.
Danh sách này được sao chép xuống bộ nhớ thiết bị, sau đó một nhân khởi tạo chạy
một lần sẽ dựng các đối tượng khối cầu (tĩnh hoặc chuyển động) ngay trên
GPU và gom chúng vào một danh sách vật thể chung. Việc cấp phát và dựng đối tượng
trực tiếp trên thiết bị cho phép sử dụng cơ chế đa hình (hàm ảo tìm giao) ngay
trong nhân kết xuất mà không phải sao chép con trỏ qua lại giữa hai không gian bộ
nhớ.

Khi dò tia, danh sách vật thể được duyệt tuyến tính: với mỗi tia, hệ thống lần
lượt kiểm tra giao với từng khối cầu và luôn giữ lại giao điểm gần nhất bằng cách
thu hẹp dần cận trên của tham số tia. Mỗi khối cầu được kiểm tra giao bằng cách
giải phương trình bậc hai giữa tia và mặt cầu như đã nêu ở
Chương~\ref{ch:ly-thuyet}, rồi điền vào bản ghi va chạm vị trí, pháp tuyến và vật
liệu tại điểm giao. Cách duyệt tuyến tính có độ phức tạp tỉ lệ với số vật thể nên
trở thành nút thắt hiệu năng khi cảnh đông đối tượng; đây chính là động lực cho
việc tích hợp cấu trúc tăng tốc không gian (BVH) nêu ở Chương~\ref{ch:ket-luan}.

\section{Nạp cảnh từ định dạng Mitsuba 3}
Để tái sử dụng các cảnh chuẩn và tách dữ liệu cảnh khỏi mã nguồn, hệ thống hiện
thực một bộ nạp phân tích tệp mô tả cảnh theo định dạng Mitsuba~3 bằng thư viện
\texttt{tinyxml2}. Bộ nạp duyệt các thẻ con của cảnh và ánh xạ chúng sang mô hình
nội bộ: thẻ cảm biến cung cấp trường nhìn để dựng camera; thẻ vật liệu được ánh
xạ theo loại (khuếch tán sang Lambertian, điện môi sang Dielectric, dẫn điện sang
Metal) và lưu lại theo định danh để tham chiếu; thẻ hình học với khối cầu cho biết
bán kính, còn tâm được lấy từ cột tịnh tiến của ma trận biến đổi $4\times4$, và
vật liệu được liên kết qua định danh đã lưu. Phiên bản hiện tại tập trung hỗ trợ
khối cầu; các loại hình học khác được bỏ qua kèm cảnh báo và là hướng mở rộng tự
nhiên trong tương lai.

\section{Tích lũy theo thời gian và hiệu chỉnh gamma}
Do mỗi khung chỉ lấy một mẫu cho mỗi điểm ảnh, ảnh của một khung đơn lẻ rất nhiễu.
Hệ thống duy trì một bộ đệm tích lũy lưu tổng màu qua các khung; màu hiển thị là
trung bình cộng của tổng này chia cho số khung đã tích lũy, đúng theo bộ ước lượng
Monte Carlo nên ảnh hội tụ dần và bớt nhiễu theo thời gian khi camera đứng yên.
Mỗi khi camera hoặc một tham số kết xuất thay đổi, host đặt lại bộ đếm khung về 1
để loại bỏ những mẫu cũ không còn hợp lệ với góc nhìn mới. Trước khi ghi ra bộ đệm
hiển thị, giá trị màu tuyến tính được hiệu chỉnh gamma (lấy căn bậc hai từng kênh)
rồi quy về số nguyên 8-bit cho mỗi kênh RGBA, đúng như lý do đã phân tích ở
Chương~\ref{ch:ly-thuyet}.

\section{Giao diện điều khiển tương tác}
Hệ thống cung cấp một bảng điều khiển dựng bằng thư viện ImGui, cho phép người
dùng theo dõi và tinh chỉnh quá trình kết xuất theo thời gian thực. Bảng điều
khiển hiển thị các chỉ số hiệu năng (số khung tích lũy, thời gian mỗi khung và số
khung hình trên giây) và cho phép chỉnh trực tiếp các tham số camera như trường
nhìn, khẩu độ, khoảng lấy nét cùng hệ số làm tối rìa ảnh. Song song, một bộ điều
khiển camera xử lý thao tác chuột để xoay quanh tâm cảnh, tịnh tiến và phóng
to/thu nhỏ. Bất kỳ thay đổi nào từ giao diện hay từ thao tác camera đều kích hoạt
đồng bộ tham số camera xuống GPU và đặt lại quá trình tích lũy, bảo đảm ảnh luôn
nhất quán với góc nhìn và cấu hình hiện tại.
